\documentclass[11pt,a4paper,sans]{moderncv}
\moderncvstyle{classic}
\moderncvcolor{blue}
\usepackage[scale=0.85]{geometry}
\usepackage[utf8]{inputenc}
\usepackage[T1]{fontenc}
\usepackage{lmodern}
\usepackage[ngerman]{babel}

% Personal Data
\name{Andreas}{Butry}
\title{System \& Software Engineer - Functional Safety, Engineering Processes \& AI Innovation}
\address{Louise-Schroeder-Str. 2}{67551 Worms}{Germany}
\phone[mobile]{+49~179~140~90~57}
\email{andreas.butry@gmail.com}
\social[linkedin]{Andreas Butry}

\photo[110pt][0pt]{profilbild}

\begin{document}
\makecvtitle

%====================================================
% Profil
%====================================================
\section{Profil}
\textbf{Systems \& Software Engineer} (8+ Jahre) mit Expertise in sicherheitskritischen Automotive-Systemen, \textbf{Softwarequalität} und \textbf{Softwareentwicklungsprozessen}. Umfassende Erfahrung über den gesamten System- und Software-Lebenszyklus hinweg: Vom Anforderungsmanagement und System-/Softwarearchitektur bis zu Implementierung, Verifikation und Freigabe nach ISO 26262 und ASPICE. Geprägt von einem ausgeprägten \textbf{Verifikations-Mindset} (u.\,a. Test-Driven Development, TDD) und Fokus auf die kontinuierliche Verbesserung von Entwicklungsprozessen, speziell durch den Einsatz von \textbf{KI-gestützten Methoden und Werkzeugen}.

\section{Weiterentwicklung: KI \& Cloud Engineering 2026}

\begin{itemize}
\item KI-gestützter Softwareentwicklungszyklus (SDLC): TDD, Code-Generierung und Review-Automatisierung
\item Generative KI, LLM-basierte Anwendungen und Prompt Engineering
\item Agentische Systeme, Multi-Agenten-Workflows und MCP (Model Context Protocol)
\item DevOps \& Cloud-Technologien: Docker, Git, Jenkins, Kubernetes und CI/CD-Pipelines
\end{itemize}
%====================================================
% Berufserfahrung
%====================================================
\section{Berufserfahrung}
\cventry{05/2025 -- heute}{Senior (Embedded) Software Design Engineer}{TES Electronics}{Frankfurt/Offenbach}{}{System- und Softwareentwicklung für Embedded Systems
\begin{itemize}
\item Konzeption und Entwicklung von APIs für cloudbasierte Services und vernetzte Home-Automation-Lösungen
\item Aufbau und Optimierung von Entwicklungsprozessen für KI-generierte Software mit Agentic Design (Claude Code)
\item Verantwortung für Projektplanung, technische Abstimmung und Koordination von Entwicklungsprojekten
\end{itemize}}
\cventry{09/2016 -- 05/2025}{Funktionsentwicklungsingenieur (Embedded Software)}{Continental}{Frankfurt}{}{Funktionsentwicklung für sicherheitskritische Bremssteuergeräte für Fahrzeuge\\newline{}
\begin{itemize}
\item \textbf{Produktentwicklung und -management für EU7-konforme Bremssysteme}
  \begin{itemize}
  \item Risikobewertung, Definition und Auswahl verschiedener Software-Kompensationsalgorithmen
  \item Technische Unterstützung für Entwickler, Testplanung und bereichsübergreifende Koordination
  \item Kundenkommunikation (z.\,B. BMW, VW) zu Anforderungen, Tests und Auslieferung von Testversionen
  \end{itemize}
\item \textbf{Neuentwicklung} eines modellbasierten (Simulink) Reifentoleranzausgleichs-Algorithmus unter Verwendung einer multidimensionalen linearen Regression (rekursiver Least-Squares-Schätzer)
\item \textbf{Systementwicklung}: Redundante Bremssteuerung für hochautomatisiertes Fahren (Zwei-ECU-Architektur)
  \begin{itemize}
  \item Entwicklung der Raddrehzahlsignal-Kommunikation zwischen ECUs über CAN-Bus
  \item Fehlererkennung und Degradationsstrategien zur Sicherstellung von funktionaler Sicherheit und Systemzuverlässigkeit
  \end{itemize}
\item \textbf{Softwarequalität \& Prozessverantwortung}: Entwicklung nach ASPICE- und ISO~26262-Standards
\item \textbf{Softwareverantwortung}: Reifentoleranzausgleich, Radsignalkette, Streckenneigungsschätzung
\end{itemize}}

\cventry{11/2014 -- 03/2015}{Praktikum (Fahrzeugintegration)}{ThyssenKrupp Presta}{Liechtenstein}{}{
\begin{itemize}
  \item Integration von Steer-by-Wire-Systemen
  \item Planung, Durchführung und Auswertung von Prüfstandversuchen
  \item Validierung eines Fahrdynamikmodells in Matlab/Simulink
\end{itemize}}

\cventry{10/2009 -- 03/2015}{Studentische Hilfskraft}{TU Darmstadt}{Darmstadt}{}{
\begin{itemize}
  \item Institut für Fahrzeugtechnik: Aufbau und Inbetriebnahme einer omnidirektionalen Fahrsimulatorplattform; Entwicklung des Echtzeit-Rechnersystems (IPG HIL, MATLAB/Simulink, C++)
  \item Institut für Produktionstechnik: Aufbau eines Prüfstands zur Analyse von Motorspindel-Wälzlagern; Steuerung über dSPACE ControlDesk und Datenerfassung inkl. Highspeed-Videoanalyse
  \item Mentoring von Erstsemestern sowie Übungsgruppenleitung in „Elektrotechnik I“
\end{itemize}}

%====================================================
% Bildung
%====================================================
\section{Bildung}
\cventry{10/2008 -- 08/2016}{M.Sc. Elektrotechnik und Automatisierungstechnik}{TU Darmstadt}{}{Abschlussnote: 1,9}{}

\section{Abschlussarbeiten / Projektarbeiten}
\cvitem{Masterarbeit}{\emph{„Analyse und Weiterentwicklung der Ansteueralgorithmen eines Fahrsimulators“}, Mercedes Benz, Stuttgart (Daimler), Abschlussnote: sehr gut}
\cvitem{Projektseminar}{\emph{„Modellbildung und Simulation des Bremsvorgangs durch ein hydraulisches Bremssystem“}, Abschlussnote: sehr gut}
\cvitem{Projektseminar}{\emph{„Entwurf und Umsetzung eines Konzepts zur Auswertung der Sensordaten eines Correvit-Sensors und zur Kompensation von Messfehlern mithilfe einer geeigneten Sensordatenfusion im Fahrzeug“}, Abschlussnote: sehr gut}
\cvitem{Bachelorarbeit}{\emph{„Implementierung eines Reglers in ein bestehendes Fahrdynamikmodell für die omnidirektionale Grundplattform eines selbstfahrenden Fahrsimulators“}, Mercedes Benz, Stuttgart (Daimler), Abschlussnote: sehr gut}

%====================================================
% Fähigkeiten
%====================================================
\section{Fähigkeiten}
\cvdoubleitem{\textbf{Sprachen}}{C, C++, C\#, Java, Python, HTML, CSS, SQL, UML, LaTeX}{\textbf{Modellierung}}{Matlab/Simulink, AUTOSAR, IBM Rhapsody, Enterprise Architect}

\cvdoubleitem{\textbf{Betriebssyst.}}{Windows, Linux, RTOS}{\textbf{Prozesse}}{ISO~26262, ASPICE, TDD / Verifikations-Mindset}

\cvdoubleitem{\textbf{Tools}}{Git, GitHub, Jira, DOORS, IPG CarMaker, Eclipse, Visual Studio, MISRA-C, Polyspace, dSPACE ControlDesk}{\textbf{KI}}{Prompt Engineering, Context Engineering, AI-Driven Software Development, Agentic Workflows, Retrieval-Augmented Generation (RAG), Model Context Protocol (MCP)}

%====================================================
% Sprachen
%====================================================
\section{Sprachen}
\cvitemwithcomment{Deutsch}{Muttersprache}{}
\cvitemwithcomment{Englisch}{verhandlungssicher}{}
\cvitemwithcomment{Spanisch}{B1}{}

%====================================================
% Interessen
%====================================================
\section{Interessen}
\cvitem{Sport}{Klettern, Surfen, Kitesurfen, Wandern}
\cvitem{Technik}{Jede Form von Technik}
\cvitem{Musik}{Musikproduktion, Soundengineering}

\clearpage
%-----       letter       ---------------------------------------------------------
% recipient data
\recipient{Company Recruitment team}{Company, Inc.\\123 somestreet\\some city}
\date{January 01, 1984}
\opening{Dear Sir or Madam,}
\closing{Yours faithfully,}
\enclosure[Attached]{curriculum vit\ae{}}          % use an optional argument to use a string other than "Enclosure", or redefine \enclname
\makelettertitle

Lorem ipsum dolor sit amet, consectetur adipiscing elit. Duis ullamcorper neque sit amet lectus facilisis sed luctus nisl iaculis. Vivamus at neque arcu, sed tempor quam. Curabitur pharetra tincidunt tincidunt. Morbi volutpat feugiat mauris, quis tempor neque vehicula volutpat. Duis tristique justo vel massa fermentum accumsan. Mauris ante elit, feugiat vestibulum tempor eget, eleifend ac ipsum. Donec scelerisque lobortis ipsum eu vestibulum. Pellentesque vel massa at felis accumsan rhoncus.

Suspendisse commodo, massa eu congue tincidunt, elit mauris pellentesque orci, cursus tempor odio nisl euismod augue. Aliquam adipiscing nibh ut odio sodales et pulvinar tortor laoreet. Mauris a accumsan ligula. Class aptent taciti sociosqu ad litora torquent per conubia nostra, per inceptos himenaeos. Suspendisse vulputate sem vehicula ipsum varius nec tempus dui dapibus. Phasellus et est urna, ut auctor erat. Sed tincidunt odio id odio aliquam mattis. Donec sapien nulla, feugiat eget adipiscing sit amet, lacinia ut dolor. Phasellus tincidunt, leo a fringilla consectetur, felis diam aliquam urna, vitae aliquet lectus orci nec velit. Vivamus dapibus varius blandit.

Duis sit amet magna ante, at sodales diam. Aenean consectetur porta risus et sagittis. Ut interdum, enim varius pellentesque tincidunt, magna libero sodales tortor, ut fermentum nunc metus a ante. Vivamus odio leo, tincidunt eu luctus ut, sollicitudin sit amet metus. Nunc sed orci lectus. Ut sodales magna sed velit volutpat sit amet pulvinar diam venenatis.

Albert Einstein discovered that $e=mc^2$ in 1905.

\[ e=\lim_{n \to \infty} \left(1+\frac{1}{n}\right)^n \]

\makeletterclosing


\end{document}